\documentclass{ifpr-pinhais}

\usepackage{amsmath}
\usepackage{lmodern}
\usepackage{graphicx}
\usepackage{latexsym}
\usepackage{amssymb}
\usepackage{placeins}
\usepackage{listings}
\usepackage{xcolor}

\graphicspath{{imagens/}}
\lstset{basicstyle=\ttfamily\footnotesize,columns=fixed,keepspaces=true,upquote=true,breaklines=true,postbreak=\mbox{$\hookrightarrow$\space},frame=single,showstringspaces=false}
\newcolumntype{P}[1]{>{\raggedright\arraybackslash}p{#1}}

% ---- Metadados do trabalho ----
\curso{Bacharelado em Ciência da Computação}
\departamento{Coordenação do Curso de Bacharelado em Ciência da Computação}
\autor{Beatriz Regina Nahas \\ João Pedro dos Santos Henrique Plinta}
\titulo{Manual 2 -- eBPF para Seguran\c{c}a}
\tipotrabalho{Manual}
\local{Pinhais}
\data{2026}

\orientador{Prof. Guilherme Werneck de Oliveira}

\preambulo{%
  Manual produzido no âmbito do projeto de extensão do minicurso
  Introdução ao eBPF, apresentado à disciplina de Práticas de Extensão V,
  do curso de Bacharelado em Ciência da Computação do \ifpr, \campus,
  sob a orientação do professor Guilherme Werneck de Oliveira.
}

\makeatletter
\hypersetup{
  pdftitle={\@title},
  pdfauthor={\@author},
  pdfsubject={\imprimirpreambulo},
  pdfcreator={LaTeX com abntex2 e ifpr-pinhais},
  hidelinks,
}
\makeatother

\setlength{\parskip}{0.1cm}

\makeindex

\begin{document}
\frenchspacing

\pretextual
\imprimircapa
\imprimirfolhaderosto{\folhaderosto}
\pdfbookmark[0]{\contentsname}{toc}
\tableofcontents*
\cleardoublepage

\textual
\pagestyle{simple}

\chapter{Introdução}

\section{eBPF no contexto de segurança}

\section{Casos de uso de segurança}

\section{Vantagens do eBPF para segurança}

\chapter{Pré-requisitos}

\section{Conhecimentos necessários}

\section{Ambiente Linux}

\chapter{Fundamentos de Segurança com eBPF}

\section{Monitoramento de eventos}

\section{Rastreamento de processos}

\section{Monitoramento de chamadas de sistema}

\section{Controle de acesso}

\section{Detecção de comportamentos suspeitos}

\chapter{eBPF e o Kernel Linux}

\section{Interação com o kernel}

\section{Hooks}

\section{Programas eBPF}

\section{Verifier}

\section{Maps}

\chapter{Mecanismos de Segurança}

O eBPF não é apenas uma ferramenta de observação passiva: integrado a subsistemas
de segurança consolidados do kernel Linux, ele viabiliza o controle ativo de
processos, redes e recursos do sistema com um nível de programabilidade que
abordagens tradicionais não oferecem \cite{linuxkernel2024}.
Este capítulo apresenta os principais mecanismos pelos quais o eBPF interage com
a infraestrutura de segurança do kernel.

% ----------------------------------------------------------------
\section{eBPF LSM}
% ----------------------------------------------------------------

O \textit{Linux Security Module} (LSM) é uma interface do kernel que define um
conjunto de \textit{hooks} em pontos de decisão de segurança — abertura de
arquivos, criação de processos, operações de rede, acesso a capacidades do sistema,
entre outros — e permite que módulos de segurança interceptem e influenciem essas
operações \cite{linuxkernel2024}.
Mecanismos tradicionais como SELinux e AppArmor são implementados como módulos
LSM, porém requerem configuração complexa e não podem ser adaptados em tempo de
execução.

A partir do kernel 5.7, o tipo de programa \texttt{BPF\_PROG\_TYPE\_LSM} permite
anexar programas eBPF diretamente aos \textit{hooks} LSM, criando módulos de
segurança totalmente programáveis \cite{linuxkernel2024}.
Ao retornar um valor de erro (por exemplo, \texttt{-EPERM}), o programa bloqueia
a operação antes de sua conclusão; ao retornar zero, permite que ela prossiga.
Essa decisão ocorre no próprio kernel, com latência de microssegundos, sem
qualquer envolvimento do espaço de usuário no caminho crítico
\cite{gregg2019}.

Entre os \textit{hooks} LSM mais relevantes para segurança estão:

\begin{itemize}
  \item \texttt{bpf\_lsm\_file\_open} — intercepta a abertura de arquivos,
        permitindo negar acesso com base em atributos do processo solicitante ou
        do arquivo;
  \item \texttt{bpf\_lsm\_bprm\_check\_security} — avalia a execução de um novo
        programa, viabilizando \textit{allowlists} de binários autorizados;
  \item \texttt{bpf\_lsm\_socket\_connect} — controla o estabelecimento de
        conexões de rede por processo, útil para isolar serviços comprometidos;
  \item \texttt{bpf\_lsm\_task\_kill} — regula o envio de sinais entre processos,
        prevenindo interferências maliciosas entre contêineres.
\end{itemize}

Uma vantagem significativa do LSM BPF em relação ao SELinux e ao AppArmor é a
possibilidade de combinar múltiplos programas no mesmo \textit{hook} — graças ao
mecanismo \textit{BPF link} — e de atualizar as políticas dinamicamente, sem
reiniciar serviços ou carregar novos módulos de kernel \cite{linuxkernel2024}.

% ----------------------------------------------------------------
\section{\textit{Cgroup hooks}}
% ----------------------------------------------------------------

\textit{Control groups} (cgroups) são o mecanismo do kernel Linux para organizar
processos em hierarquias e aplicar limites de recursos (CPU, memória, I/O de
rede e disco) a cada grupo \cite{linuxkernel2024}.
Em ambientes de contêineres, cada contêiner é tipicamente mapeado a um cgroup
distinto, tornando os \textit{hooks} de cgroup um ponto natural de aplicação de
políticas de segurança com escopo por contêiner \cite{gregg2020}.

O eBPF oferece um conjunto de tipos de programa associados a cgroups que permitem
interceptar operações antes que elas aconteçam:

\textbf{\texttt{BPF\_PROG\_TYPE\_CGROUP\_SKB}.}
Filtra pacotes na entrada ou saída de um cgroup específico.
Permite implementar políticas de rede por contêiner sem depender de regras de
\textit{iptables} ou \textit{nftables}, com menor overhead e maior expressividade
\cite{linuxkernel2024}.

\textbf{\texttt{BPF\_PROG\_TYPE\_CGROUP\_SOCK} e variantes.}
Controla a criação e o uso de sockets por processos do cgroup.
Pode bloquear a criação de sockets de determinados protocolos ou famílias de
endereço, restringindo efetivamente a superfície de ataque de rede de um
contêiner \cite{gregg2019}.

\textbf{\texttt{BPF\_PROG\_TYPE\_CGROUP\_DEVICE}.}
Regula o acesso a dispositivos (\texttt{/dev/*}) por processos do cgroup,
substituindo com vantagens a lista de dispositivos do cgroups v1
\cite{linuxkernel2024}.

\textbf{\texttt{BPF\_PROG\_TYPE\_CGROUP\_SYSCTL}.}
Intercepta leituras e escritas em parâmetros \texttt{sysctl} realizadas por
processos do cgroup, impedindo que um contêiner altere configurações globais do
kernel que afetariam outros contêineres \cite{linuxkernel2024}.

A granularidade por cgroup torna esses \textit{hooks} particularmente adequados
para ambientes Kubernetes, onde políticas de segurança precisam ser aplicadas
individualmente a cada \textit{pod} sem interromper os demais
\cite{cassagnes2020}.

% ----------------------------------------------------------------
\section{Seccomp e eBPF}
% ----------------------------------------------------------------

O \textit{secure computing mode} (seccomp) é um mecanismo do kernel Linux que
restringe o conjunto de chamadas de sistema que um processo pode invocar
\cite{linuxkernel2024}.
No modo \texttt{SECCOMP\_MODE\_FILTER}, introduzido no kernel 3.5, um programa BPF
clássico (cBPF) define a política de filtragem — especificando quais chamadas de
sistema são permitidas, bloqueadas ou encaminhadas para análise em espaço de usuário
\cite{mccanne1993}.

Com o eBPF, dois aspectos dessa integração evoluíram significativamente:

\textbf{Geração programática de filtros seccomp.}
Ferramentas modernas como \texttt{libseccomp} compilam políticas expressas em alto
nível para o formato cBPF exigido pelo kernel.
Pesquisas recentes exploram o uso de eBPF para gerar esses filtros de forma
dinâmica a partir de perfis de comportamento observados em execução — em vez de
listas estáticas definidas em tempo de desenvolvimento \cite{gregg2019}.

\textbf{Auditoria de violações seccomp com eBPF.}
Quando um filtro seccomp retorna \texttt{SECCOMP\_RET\_TRACE} ou
\texttt{SECCOMP\_RET\_LOG}, a chamada de sistema bloqueada pode ser capturada por
um programa eBPF conectado ao \textit{tracepoint} \texttt{syscalls:sys\_enter\_*},
permitindo registrar tentativas de evasão com contexto completo — processo, PID,
UID, argumentos — sem modificar o filtro seccomp original \cite{linuxkernel2024}.

A combinação de seccomp com LSM BPF representa a abordagem de defesa em
profundidade (\textit{defense in depth}) recomendada para ambientes de produção:
o seccomp restringe o conjunto de chamadas de sistema acessíveis, enquanto o LSM
BPF impõe controles adicionais sobre o que pode ser feito com as chamadas
permitidas \cite{beyer2016}.
Juntos, os dois mecanismos reduzem drasticamente a superfície de ataque exposta
por um processo comprometido \cite{gregg2020}.

% ----------------------------------------------------------------
\section{Controle de processos e acesso a recursos}
% ----------------------------------------------------------------

Além dos mecanismos anteriores — orientados a interfaces específicas como sockets,
arquivos e chamadas de sistema —, o eBPF oferece visibilidade e controle sobre o
ciclo de vida de processos e o uso de recursos do sistema operacional de forma
holística \cite{gregg2019}.

\textbf{Rastreamento de ancestralidade de processos.}
Utilizando mapas BPF do tipo \texttt{BPF\_MAP\_TYPE\_HASH}, é possível manter em
tempo real uma tabela de processos ativos, registrando PID, PPID, UID, executável
e cgroup de cada processo \cite{linuxkernel2024}.
Essa árvore de processos viabiliza a detecção de \textit{process injection} — em
que um atacante insere código em um processo legítimo —, pois alterações inesperadas
na identidade ou no comportamento de um processo tornam-se visíveis
\cite{majors2022}.

\textbf{Monitoramento de uso de capacidades (\textit{capabilities}).}
O modelo de capacidades do Linux substitui o binômio root/não-root por um conjunto
granular de privilégios (\texttt{CAP\_NET\_ADMIN}, \texttt{CAP\_SYS\_PTRACE},
\texttt{CAP\_DAC\_OVERRIDE}, entre outras) \cite{linuxkernel2024}.
Programas eBPF conectados ao \textit{hook} LSM \texttt{bpf\_lsm\_capable} podem
registrar toda tentativa de uso de capacidade, identificando processos que operam
com mais privilégios do que o esperado — um indicador comum de escalonamento de
privilégios \cite{gregg2020}.

\textbf{Controle de acesso a namespaces.}
Namespaces do Linux isolam recursos como PID, rede, sistema de arquivos e
identificadores de usuário entre grupos de processos \cite{linuxkernel2024}.
Tentativas de um processo sair do seu namespace — por exemplo, via
\texttt{setns(2)} ou \texttt{unshare(2)} — podem ser monitoradas e bloqueadas
por programas eBPF, prevenindo fugas de contêineres (\textit{container escapes})
\cite{cassagnes2020}.

\textbf{Observabilidade de acesso a arquivos e descritores.}
O \textit{tracepoint} \texttt{sched\_process\_exit}, combinado com um mapa de
estado de arquivos abertos, permite detectar processos que terminam sem fechar
descritores sensíveis — padrão observado em algumas classes de \textit{malware}
que abusam de herança de descritores entre processos \cite{gregg2019}.

\textbf{Limitação dinâmica de recursos.}
Em conjunto com os \textit{cgroup hooks} descritos na seção anterior, programas
eBPF podem implementar cotas dinâmicas de chamadas de sistema por processo — por
exemplo, limitando a taxa de chamadas \texttt{fork(2)} para mitigar ataques de
exaustão de processos (\textit{fork bombs}) — sem necessidade de reconfigurar
os limites de cgroup em espaço de usuário \cite{linuxkernel2024}.

A combinação desses controles com os mecanismos de observabilidade descritos no
capítulo anterior cria uma camada de segurança coesa: eventos suspeitos são
detectados, correlacionados e, quando necessário, bloqueados em uma única
infraestrutura baseada em eBPF, sem a sobrecarga de agentes externos e com
latência compatível com ambientes de produção \cite{majors2022}.
\chapter{Detecção e Monitoramento}
\label{cap:deteccao}

A capacidade de detectar atividades maliciosas em tempo real, sem degradar o
desempenho das aplicações monitoradas, é um dos principais diferenciais do eBPF
frente a soluções tradicionais de segurança \cite{cassagnes2020}.
Este capítulo apresenta técnicas concretas de detecção e monitoramento baseadas
em eBPF, organizadas pelos vetores de observação mais relevantes em ambientes de
produção.

% ----------------------------------------------------------------
\section{Detecção de execução de processos}
% ----------------------------------------------------------------

A execução de processos é um dos eventos mais informativos para segurança:
praticamente toda ação maliciosa — desde a exploração de uma vulnerabilidade até
a persistência no sistema — envolve em algum momento a criação ou substituição de
um processo \cite{gregg2019}.
O eBPF permite monitorar esse vetor com cobertura total e latência mínima,
independentemente da origem do processo.

O \textit{tracepoint} \texttt{sys\_enter\_execve} é acionado toda vez que um
processo invoca a chamada de sistema \texttt{execve(2)}, responsável por substituir
a imagem do processo atual por um novo executável \cite{linuxkernel2024}.
Um programa eBPF nesse ponto coleta, sem modificar o processo:

\begin{itemize}
  \item PID e PPID — identificadores do processo e de seu pai, permitindo
        reconstruir a árvore de execução;
  \item UID e GID — identidade real do solicitante, revelando tentativas de
        execução sob identidade indevida;
  \item caminho do executável e argumentos da linha de comando, via
        \texttt{bpf\_probe\_read\_user\_str};
  \item cgroup e namespace de PID, essenciais para identificar o contêiner de
        origem em ambientes Kubernetes \cite{cassagnes2020}.
\end{itemize}

O exemplo a seguir ilustra um script bpftrace que registra toda execução
originada de processos filhos de servidores web — padrão típico de exploração de
vulnerabilidades de execução remota de código:

\begin{lstlisting}[language=bash, caption={Detecção de shell a partir de processo de serviço web com bpftrace.}]
tracepoint:syscalls:sys_enter_execve
/comm == "nginx" || comm == "apache2"/
{
    printf("[ALERTA] shell filho de %s: %s (pid=%d uid=%d)\n",
           comm, str(args->filename), pid, uid);
}
\end{lstlisting}

Além do \texttt{execve}, o \textit{tracepoint} \texttt{sched\_process\_fork}
captura a criação de processos via \texttt{fork(2)} e \texttt{clone(2)},
permitindo rastrear a propagação de processos maliciosos que se replicam antes de
executar sua carga \cite{gregg2020}.
Combinados em um mapa BPF que mantém a árvore de parentesco em memória, esses
dois pontos de observação fornecem uma visão completa do ciclo de vida de processos
sem qualquer instrumentação das aplicações \cite{majors2022}.

% ----------------------------------------------------------------
\section{Monitoramento de arquivos, chamadas de sistema e conexões}
% ----------------------------------------------------------------

\textbf{Acesso a arquivos.}
Programas eBPF conectados ao \textit{tracepoint} \texttt{sys\_enter\_openat}
registram toda abertura de arquivo com contexto completo — processo, caminho,
\textit{flags} de acesso e resultado da operação \cite{linuxkernel2024}.
Caminhos de particular interesse para segurança incluem:

\begin{itemize}
  \item \texttt{/etc/shadow}, \texttt{/etc/passwd} e \texttt{/root/.ssh/} —
        arquivos de credenciais frequentemente alvo de exfiltração;
  \item \texttt{/proc/\$PID/mem} e \texttt{/proc/\$PID/maps} — interfaces de
        memória de processo usadas por técnicas de injeção \cite{manpages2024};
  \item \texttt{/usr/bin/}, \texttt{/lib/} e \texttt{/sbin/} abertos em modo de
        escrita — indicativo de \textit{rootkits} que modificam binários do sistema.
\end{itemize}

A ferramenta \texttt{opensnoop} do BCC \cite{bcc} implementa esse padrão e é
amplamente utilizada em investigações forenses para identificar quais arquivos
foram acessados durante um incidente.

\textbf{Chamadas de sistema.}
O conjunto completo de \textit{tracepoints} \texttt{sys\_enter\_*} oferece
visibilidade sobre qualquer chamada de sistema invocada por qualquer processo
\cite{linuxkernel2024}.
Para fins de detecção, chamadas críticas são monitoradas seletivamente para
reduzir o volume de eventos:

\begin{itemize}
  \item \texttt{ptrace} — interceptação de outro processo; abusada por
        \textit{malware} para injeção de código e por \textit{rootkits} para
        ocultar presença;
  \item \texttt{setuid} / \texttt{setgid} / \texttt{setcap} — alteração de
        identidade ou capacidades, indicativas de escalonamento de privilégios;
  \item \texttt{mmap} com \texttt{PROT\_EXEC} — mapeamento de memória executável,
        padrão de carregamento de shellcode \cite{gregg2019};
  \item \texttt{init\_module} / \texttt{finit\_module} — carregamento de módulos
        de kernel, vetor clássico de instalação de \textit{rootkits}.
\end{itemize}

\textbf{Conexões de rede.}
Combinando \textit{tracepoints} de \textit{socket} com \textit{hooks} de XDP e
TC, o eBPF oferece visibilidade em múltiplas camadas da pilha de rede
\cite{linuxkernel2024}.
O monitoramento de conexões de saída — via \texttt{sys\_enter\_connect} — permite
identificar processos que estabelecem comunicação com endereços não previstos,
padrão associado a \textit{beaconing} de \textit{Command and Control} (C2) e
exfiltração de dados \cite{cassagnes2020}.
A ferramenta \texttt{tcpconnect} do BCC \cite{bcc} registra cada nova conexão TCP
com PID, endereço de destino e porta, sem impacto perceptível sobre o throughput
da rede \cite{gregg2019}.

% ----------------------------------------------------------------
\section{Detecção de comportamento anômalo}
% ----------------------------------------------------------------

A detecção de anomalias vai além do monitoramento de eventos individuais: consiste
em identificar sequências ou padrões de comportamento que diferem do que é esperado
para um processo, serviço ou usuário \cite{majors2022}.
O eBPF é particularmente adequado para esse fim porque permite acumular estado em
mapas BPF ao longo do tempo e realizar correlações diretamente no kernel, antes
mesmo que os eventos sejam enviados ao espaço de usuário \cite{linuxkernel2024}.

\textbf{Desvio de perfil de chamadas de sistema.}
Cada classe de aplicação — servidor web, banco de dados, agente de monitoramento
— apresenta um perfil estável de chamadas de sistema em operação normal.
Um programa eBPF pode construir, durante um período de observação inicial, um
histograma de frequência de chamadas por processo e alertar quando a distribuição
observada diverge significativamente desse perfil de referência
\cite{gregg2019}.
Essa técnica é resistente a \textit{malware} que tenta mimetizar nomes de
processos legítimos, pois o comportamento em nível de chamadas de sistema é mais
difícil de falsificar.

\textbf{Comunicação lateral entre contêineres.}
Em ambientes Kubernetes, comunicação direta entre \textit{pods} que normalmente
não se comunicam é um indicador de movimentação lateral após comprometimento
\cite{cassagnes2020}.
Mapas BPF do tipo \texttt{BPF\_MAP\_TYPE\_HASH} indexados por par
(cgroup de origem, cgroup de destino) permitem detectar esses fluxos anômalos em
tempo real, sem depender de logs de \textit{firewall} com atraso de segundos
\cite{linuxkernel2024}.

\textbf{Escalada de privilégios em tempo de execução.}
A sequência \texttt{fork} → \texttt{execve} de um binário SUID → \texttt{setuid(0)}
é o padrão de exploração de vulnerabilidades de escalonamento de privilégios local
mais comum \cite{gregg2020}.
Um programa eBPF que rastreia transições de UID em mapas indexados por PID detecta
essa sequência em microssegundos, muito antes que ela se manifeste em logs de
auditoria \cite{gregg2019}.

\textbf{Injeção de código em processos em execução.}
O uso de \texttt{ptrace} ou de \texttt{process\_vm\_writev} para escrever em
memória de outro processo é uma técnica de injeção amplamente explorada
\cite{gregg2019}.
Correlacionando o PID do processo chamador com o PID alvo e verificando se a
relação pai-filho é legítima — usando a árvore de processos mantida em um mapa
BPF — é possível distinguir operações legítimas de depuradores de tentativas de
injeção maliciosa \cite{majors2022}.

% ----------------------------------------------------------------
\section{Auditoria de eventos}
% ----------------------------------------------------------------

A auditoria de segurança tem como objetivo registrar, de forma confiável e
verificável, eventos relevantes para conformidade e investigação forense
\cite{beyer2016}.
O subsistema \texttt{auditd} do Linux é a solução tradicional para esse fim, mas
apresenta limitações de desempenho em sistemas de alta carga e dificuldades de
extensão para novos tipos de eventos \cite{gregg2020}.
O eBPF surge como complemento ou substituto natural: coleta eventos com menor
sobrecarga, permite filtragem inteligente diretamente no kernel e é extensível sem
alterações no sistema operacional \cite{cassagnes2020}.

\textbf{Registro de eventos com contexto completo.}
Por meio do mapa \texttt{BPF\_MAP\_TYPE\_RINGBUF}, programas eBPF emitem registros
de auditoria estruturados para o espaço de usuário com campos como timestamp em
nanossegundos, PID, UID, cgroup, chamada de sistema, argumentos e código de retorno
\cite{linuxkernel2024}.
Essa riqueza de contexto supera a granularidade oferecida pelo \texttt{auditd}
para a maioria dos eventos \cite{gregg2019}.

\textbf{Filtragem inteligente no kernel.}
Diferentemente do \texttt{auditd}, que envia todos os eventos para o espaço de
usuário e filtra depois, programas eBPF aplicam regras de interesse diretamente
no kernel — registrando apenas eventos que satisfazem critérios de segurança
predefinidos \cite{linuxkernel2024}.
Isso reduz drasticamente o volume de dados transmitidos e processados, tornando
viável a auditoria contínua mesmo em sistemas com dezenas de milhares de chamadas
de sistema por segundo \cite{cassagnes2020}.

\textbf{Integridade do registro de auditoria.}
Uma preocupação legítima em auditoria forense é a possibilidade de um atacante com
privilégios suficientes suprimir registros de suas ações.
Programas eBPF em execução no kernel são mais difíceis de desabilitar
silenciosamente do que agentes em espaço de usuário: sua remoção requer privilégios
explícitos e gera eventos observáveis por outros programas eBPF — criando
uma cadeia de custódia parcialmente auto-referencial \cite{majors2022}.

\textbf{Integração com sistemas externos.}
Registros de auditoria coletados por programas eBPF podem ser encaminhados para
sistemas de gerenciamento de eventos e informações de segurança (SIEM) via
exportadores como o \texttt{ebpf\_exporter} \cite{ebpfexporter} ou via
\textit{pipelines} Prometheus/Grafana \cite{prometheusdocs,grafanadocs}.
Essa integração centraliza os dados de múltiplos nós em uma visão unificada,
essencial para detecção de ataques distribuídos que envolvem vários sistemas
simultaneamente \cite{beyer2016}.
\chapter{Ferramentas de Segurança}

\section{BCC}

\section{bpftrace}

\section{libbpf}

\section{Falco}

\section{Tetragon}

\chapter{eBPF em Containers e Kubernetes}
\label{cap:containers}

A adoção massiva de contêineres e orquestradores como o Kubernetes trouxe novos
desafios de segurança: ambientes altamente dinâmicos, com processos efêmeros,
múltiplos inquilinos compartilhando o mesmo kernel e superfície de ataque
distribuída por dezenas ou centenas de nós \cite{cassagnes2020}.
Abordagens tradicionais de segurança baseadas em agentes por processo ou em
inspeção de logs com atraso mostram-se inadequadas nesse contexto.
O eBPF, por operar no kernel compartilhado por todos os contêineres de um nó,
oferece visibilidade e controle centralizados sem depender de instrumentação
individual de cada carga de trabalho \cite{linuxkernel2024}.

% ----------------------------------------------------------------
\section{Segurança de contêineres}
% ----------------------------------------------------------------

Contêineres não são máquinas virtuais: compartilham o kernel do sistema hospedeiro
e dependem de mecanismos de isolamento do próprio Linux — namespaces, cgroups e
capacidades — para separar cargas de trabalho \cite{linuxkernel2024}.
Essa arquitetura torna o kernel o ponto central de controle e, consequentemente,
o principal vetor de ataque em ambientes de contêineres \cite{gregg2020}.

As ameaças mais relevantes nesse contexto são:

\textbf{Fuga de contêiner (\textit{container escape}).}
Um processo dentro de um contêiner que consegue sair do seu namespace de PID,
de rede ou de sistema de arquivos e interagir com o kernel ou com outros contêineres
sem restrição \cite{linuxkernel2024}.
Fugas exploram tipicamente vulnerabilidades em chamadas de sistema como
\texttt{unshare(2)}, \texttt{setns(2)} ou \texttt{mount(2)}, ou abusam de
capacidades Linux concedidas desnecessariamente ao contêiner.

\textbf{Escalada de privilégios dentro do contêiner.}
Mesmo sem sair do namespace, um processo que obtém \texttt{CAP\_SYS\_ADMIN}
ou executa como \texttt{root} no contêiner pode comprometer outros processos do
mesmo nó \cite{gregg2019}.

\textbf{Ataques de movimentação lateral.}
Um contêiner comprometido que estabelece comunicação direta com outros contêineres
ou serviços de infraestrutura — banco de dados, API do Kubernetes — fora dos
fluxos previstos \cite{cassagnes2020}.

O eBPF endereça esses vetores de múltiplas formas.
Programas conectados a \textit{tracepoints} de chamadas de sistema como
\texttt{unshare}, \texttt{setns} e \texttt{clone} detectam tentativas de
manipulação de namespaces em tempo real \cite{linuxkernel2024}.
\textit{Hooks} LSM BPF podem bloquear essas operações para processos cujos cgroups
não estejam em uma lista de permissões, sem qualquer modificação nas imagens de
contêiner \cite{gregg2019}.
Mapas BPF indexados por par de cgroups registram todo fluxo de rede entre
contêineres, tornando a movimentação lateral imediatamente visível \cite{majors2022}.

% ----------------------------------------------------------------
\section{Segurança em Kubernetes}
% ----------------------------------------------------------------

O Kubernetes introduz uma camada adicional de complexidade: além dos processos
em execução nos nós, existe o plano de controle — API Server, etcd, scheduler,
controller manager — cujo comprometimento implica controle total sobre o cluster
\cite{beyer2016}.
O eBPF é utilizado em duas frentes complementares nesse ambiente.

\textbf{Substituição do kube-proxy.}
O componente tradicional \texttt{kube-proxy} implementa as regras de roteamento
de serviços Kubernetes por meio de tabelas \textit{iptables} de crescimento
linear com o número de serviços \cite{linuxkernel2024}.
Projetos como o Cilium substituem o \texttt{kube-proxy} por programas eBPF anexados
a \textit{hooks} TC e XDP, implementando o mesmo roteamento com complexidade
constante e menor latência, além de oferecer visibilidade sobre cada fluxo de
rede entre \textit{pods} \cite{cassagnes2020}.

\textbf{Observabilidade do plano de dados.}
Em um cluster com centenas de \textit{pods} espalhados por dezenas de nós, rastrear
um fluxo de requisição de ponta a ponta requer correlação de eventos de múltiplos
processos e nós \cite{majors2022}.
Programas eBPF nos nós coletam identificadores de fluxo — endereços IP, portas,
metadados de cgroup — e os associam a metadados Kubernetes (namespace, nome do
\textit{pod}, \textit{labels}) obtidos do API Server, compondo uma visão de
observabilidade de rede enriquecida sem necessidade de \textit{sidecars} ou
modificação das aplicações \cite{linuxkernel2024}.

\textbf{Auditoria do API Server.}
Toda operação no cluster — criação de \textit{pods}, modificação de segredos,
concessão de permissões RBAC — passa pelo API Server e pode ser auditada por
programas eBPF conectados aos \textit{tracepoints} de rede do nó onde ele executa,
complementando os logs de auditoria nativos do Kubernetes com contexto de processo
e identidade de rede \cite{gregg2019}.

% ----------------------------------------------------------------
\section{Políticas de segurança}
% ----------------------------------------------------------------

O Kubernetes oferece primitivos nativos de política de segurança —
\textit{NetworkPolicies}, \textit{PodSecurityAdmission} e RBAC — mas sua
implementação e granularidade dependem dos plugins de infraestrutura instalados
\cite{linuxkernel2024}.
O eBPF amplia significativamente o espaço de políticas aplicáveis.

\textbf{Políticas de rede por identidade.}
Enquanto \textit{NetworkPolicies} nativas operam sobre endereços IP e portas, que
mudam a cada recriação de \textit{pod}, soluções baseadas em eBPF como o Cilium
implementam políticas baseadas em identidade: a identidade de um \textit{pod} é
derivada dos seus \textit{labels} Kubernetes e codificada em um rótulo numérico
propagado em cabeçalhos de pacotes \cite{cassagnes2020}.
Isso permite expressar regras como ``\textit{pods} do serviço A só podem receber
tráfego de \textit{pods} do serviço B no mesmo namespace'' de forma estável,
independentemente de endereços IP efêmeros \cite{linuxkernel2024}.

\textbf{Políticas de chamadas de sistema por \textit{pod}.}
Perfis seccomp restringem as chamadas de sistema disponíveis para um contêiner,
mas são estáticos e difíceis de calibrar sem causar falsos positivos \cite{gregg2020}.
O eBPF permite gerar perfis dinâmicos observando o comportamento real de cada
\textit{pod} durante um período de aprendizado e, em seguida, aplicar uma política
derivada desse perfil sem modificar o manifesto do \textit{pod}
\cite{nakryiko2020}.

\textbf{\textit{TracingPolicies} declarativas.}
Ferramentas como o Tetragon expõem um recurso Kubernetes
(\texttt{TracingPolicy}) que permite expressar políticas de segurança em YAML —
quais chamadas de sistema monitorar, quais argumentos filtrar e qual ação tomar
(alertar, registrar ou bloquear) — que são traduzidas automaticamente em programas
eBPF e implantadas em todos os nós do cluster \cite{linuxkernel2024}.
Esse modelo declarativo integra a segurança ao fluxo de \textit{GitOps} já adotado
por equipes de plataforma, tratando políticas de segurança como código versionável
e revisável \cite{beyer2016}.

% ----------------------------------------------------------------
\section{Detecção de atividades suspeitas}
% ----------------------------------------------------------------

A natureza dinâmica do Kubernetes — onde \textit{pods} são criados e destruídos
em segundos — exige que a detecção de atividades suspeitas opere em tempo real e
seja capaz de correlacionar eventos efêmeros antes que evidências desapareçam
\cite{majors2022}.
O eBPF é a tecnologia que melhor satisfaz esse requisito, e os padrões de ataque
mais comuns em ambientes Kubernetes têm contrapartidas diretas em programas de
detecção baseados em eBPF.

\textbf{Execução de shell em contêiner.}
A abertura de um terminal interativo (\texttt{/bin/sh}, \texttt{/bin/bash})
dentro de um contêiner de produção é um forte indicador de comprometimento ou
de acesso não autorizado via \texttt{kubectl exec} \cite{cassagnes2020}.
Programas eBPF detectam esse padrão combinando o \textit{tracepoint}
\texttt{sys\_enter\_execve} com a verificação de que o processo possui um TTY
associado e que seu cgroup corresponde a um \textit{pod} em execução
\cite{linuxkernel2024}.

\textbf{Acesso a segredos montados.}
Segredos Kubernetes são montados como arquivos no sistema de arquivos do contêiner.
Leitura desses arquivos por processos diferentes do esperado — especialmente por
shells ou binários de sistema — é detectada por programas eBPF conectados ao
\textit{tracepoint} \texttt{sys\_enter\_openat}, filtrando caminhos que
correspondam a diretórios de montagem de segredos (\texttt{/var/run/secrets/})
\cite{gregg2019}.

\textbf{Comunicação com o API Server a partir de contêineres de aplicação.}
Requisições HTTP ao endereço do API Server (\texttt{kubernetes.default.svc})
originadas de \textit{pods} de aplicação que não deveriam ter acesso ao plano
de controle indicam possível tentativa de movimentação para comprometer o cluster
\cite{beyer2016}.
Mapas BPF indexados por endereço de destino e cgroup de origem permitem detectar
esses fluxos em microssegundos e bloqueá-los via \textit{hook} de cgroup de socket
\cite{linuxkernel2024}.

\textbf{Modificação de imagens em tempo de execução.}
Contêineres devem ser imutáveis: qualquer escrita no sistema de arquivos raiz
fora de volumes montados explicitamente é suspeita \cite{gregg2020}.
Programas eBPF conectados a \textit{hooks} LSM de abertura de arquivo em modo de
escrita — filtrados pelo namespace de montagem do contêiner — detectam esse padrão
independentemente da ferramenta ou técnica usada pelo atacante para realizar a
escrita \cite{linuxkernel2024}.

\textbf{Varredura de rede interna.}
Um \textit{pod} comprometido que tenta conectar a múltiplos endereços e portas
em sequência rápida executa varredura de rede para mapear serviços internos do
cluster \cite{cassagnes2020}.
Mapas BPF do tipo \texttt{BPF\_MAP\_TYPE\_LRU\_HASH} acumulam contagens de
tentativas de conexão por \textit{pod} em janelas de tempo deslizantes e acionam
alertas quando o número de destinos distintos excede um limiar configurável
\cite{linuxkernel2024,gregg2020}.
\chapter{Exemplos Práticos}

\section{Monitoramento de processos e arquivos}

\section{Rastreamento de chamadas de sistema}

\section{Detecção de conexões suspeitas}

\chapter{Segurança do próprio eBPF}

\section{Privilégios necessários}

\section{eBPF Verifier}

\section{Riscos de segurança}

\section{Limitações}

\chapter{Conclusão}
\label{cap:conclusao-seguranca}

Este manual percorreu o eBPF sob a ótica da segurança. Partimos dos fundamentos
(Capítulos~\ref{cap:introducao-seguranca} a~\ref{cap:kernel-seguranca}),
passamos pelos mecanismos que o kernel oferece para controle ativo, como LSM,
\textit{cgroup hooks} e seccomp (Capítulo~\ref{cap:mecanismos}), e pelas técnicas
de detecção e auditoria (Capítulo~\ref{cap:deteccao}). Em seguida, vimos as
ferramentas que materializam essas técnicas (Capítulo~\ref{cap:ferramentas-seguranca})
e como elas se estendem a contêineres e Kubernetes (Capítulo~\ref{cap:containers}).

O Capítulo~\ref{cap:exemplos-seguranca} reuniu essas peças em três cenários de
investigação, e o Capítulo~\ref{cap:seguranca-ebpf} lembrou que o eBPF também é
superfície de ataque: os mesmos privilégios que o tornam poderoso exigem
\texttt{CAP\_BPF} restrito, o \textit{verifier} ativo e atenção às limitações de
cada ferramenta. Detectar e bloquear com eBPF só é confiável quando o próprio
mecanismo é protegido.

Os comandos do Capítulo~\ref{cap:exemplos-seguranca} estão reunidos no script
\texttt{testes/08-seguranca-exemplos.sh}, na pasta \texttt{testes/} do repositório,
para serem executados em uma máquina de laboratório com \texttt{sudo}. Quem for
adiante pode combinar o que aprendeu aqui com o Manual~1, de observabilidade,
para enviar os eventos de segurança a um painel, ou ler a documentação do Falco e
do Tetragon para levar essas regras a um ambiente real.


\postextual

% Descomente quando houver \cite{} no texto:
\bibliography{referencias}

\end{document}
